% Transformation of a VECTOR field under a rotation R, via a bar magnet whose
% magnetization m(x) is the field. A vector field is a set of ARROWS AT POSITIONS,
% and a rotation acts on BOTH independently: it moves each arrow's basepoint
% (positions, "r") AND turns each arrow (vectors, "m"). A correct transformation
% does both together, (R.m)(x) = R m(R^-1 x); doing only one leaves the arrows
% attached to points that no longer line up with the direction they point.
%
% Commutative diagram (R acts on two different spaces, so the two partial
% rotations are genuine intermediate objects that both complete to the correct
% result):
%
%   m ---- rotate r ----> (positions rotated, vectors not)   [TL -> TR]
%   |                              |
%  rotate m                      rotate m
%   v                              v
%   (vectors rotated) - rotate r -> (both) = correct         [BL -> BR]
%
% plus the direct diagonal TL -> BR. Every path flows into the correct result.
% See _template.tex / _preamble.tex for the shared setup.
% alt: A bar magnet with magnetization m and position vector x, and three
%   attempts to rotate it by R: rotating both x and m gives the correct
%   transformed field; rotating only m or only x each leaves the magnetization
%   misaligned with the bar and is marked wrong.
% width: 560
\ifnum\pdfoutput>0
\documentclass[tikz,border=3pt]{standalone}         % pdflatex → PDF proof
\else
\documentclass[tikz,border=3pt,dvisvgm]{standalone} % latex   → web SVG
\fi
% Shared setup for all site figures — \input by every figures/*.tex.
% Change the palette, font size, or driver logic HERE and every figure inherits
% it on next rebuild. See src/assets/blog/README.md for the full workflow.
%
% NOTE: this file is \input AFTER \documentclass (it needs amsmath/tikz loaded),
% but the \documentclass line itself carries the driver switch and must stay in
% each figure — see _template.tex. Keep figure .tex files starting from that
% template so the driver logic is never forgotten.

\usepackage{amsmath}

% --- Palette: sentinel hexes = the LIGHT-theme values from DESIGN.md. ----------
% The PDF proof / Overleaf output therefore look like the light rendering, and
% fig.sh rewrites each hex to its themed CSS variable in the web SVG:
%   ink    #000000 -> var(--color-text-main, currentColor)   (also all black strokes)
%   accent #003366 -> var(--color-accent, #003366)
%   muted  #555555 -> var(--color-text-muted, #555555)
%   figred #b91c1c -> var(--fig-red, #b91c1c)
%   figblue   #1d4ed8 -> var(--fig-blue, #1d4ed8)
%   figorange #c2410c -> var(--fig-orange, #c2410c)
\definecolor{ink}{HTML}{000000}
\definecolor{accent}{HTML}{003366}
\definecolor{muted}{HTML}{555555}
\definecolor{figred}{HTML}{B91C1C}
\definecolor{figblue}{HTML}{1D4ED8}
\definecolor{figorange}{HTML}{C2410C}

% --- Figure body font size = site body text. ---------------------------------
% The site renders body text at 18px ≈ 13.5pt. dvisvgm emits the SVG at its
% intrinsic point size and <Figure> shows it ~1:1, so to make figure labels match
% surrounding prose we set the figure's \normalsize to 13.5/16pt. Individual
% labels can still deviate with font=\small / font=\large as needed. Computer
% Modern matches the KaTeX body math, so $dg_p(v)$ reads identically in both.
\newcommand{\figfontsize}{\fontsize{13.5}{16}\selectfont}

% Apply it as the default font for every tikzpicture, plus sane global defaults.
% Dimension figures in mm (x=1mm,y=1mm) so geometry and this pt-based text stay
% independent — scale the drawing by editing coordinates, not the type size.
\tikzset{
  every picture/.style = {font=\figfontsize, line join=round},
  edge/.style          = {line width=1pt},
  hair/.style          = {line width=1pt},
}

\usetikzlibrary{arrows.meta}

\begin{document}
\begin{tikzpicture}[x=1mm, y=1mm]
  % --- Geometry ---------------------------------------------------------------
  \def\bw{34}   % bar length (long axis)
  \def\bh{15}   % bar height (short axis)

  \tikzset{
    mvec/.style  = {line width=1.4pt, figred, -{Stealth[length=2.2mm, width=2mm, round]}, line cap=round},
    rvec/.style  = {edge, ink, -{Stealth[length=2.4mm, width=1.9mm, round]}, line cap=round},
    tform/.style = {edge, muted, -{Stealth[length=2.4mm, width=1.9mm, round]}, line cap=round},
  }

  % \magnet{cx}{cy}{posang}{vecang} draws a magnet centered at (cx,cy).
  % POSITIONS (the body + the basepoints the field arrows sit on) are rotated by
  % posang; the arrow DIRECTIONS are rotated by vecang. These are INDEPENDENT --
  % a vector field is arrows-at-points, and "rotate r" moves the points while
  % "rotate m" turns the arrows. When posang == vecang the arrows lie along the
  % bar (physically consistent); when they differ (the wrong panels) the arrows
  % sit on the rotated lattice but point the wrong way, so field and body clash.
  \def\np{3}          % field arrows per ROW along the bar
  % Half-length of each field arrow. Basepoints are \bw/(\np+1) apart; with \np=3
  % that's ~8.5mm, so keep 2*\alen a couple mm below it for a clear gap between
  % colinear arrows.
  \def\alen{3.4}
  \def\nrow{2}        % number of arrow ROWS, symmetric about the bar's center
  \def\rowsep{6}      % vertical spacing between the two rows
  \newcommand{\magnet}[4]{%
    % Body + big faint N/S watermark live in the POSITION frame (posang).
    % `transform shape` makes node CONTENTS rotate with the scope -- without it a
    % \node's position rotates but its text stays upright, so the letters wouldn't
    % tilt with the magnet. The N/S are large, centered in each pole half, at low
    % opacity so they read as a faint label behind the field arrows (drawn after).
    \begin{scope}[shift={(#1,#2)}, rotate=#3, transform shape]
      \filldraw[edge, ink, fill=figbg] (-\bw/2,-\bh/2) rectangle (\bw/2,\bh/2);
      \draw[muted, line width=1pt] (0,-\bh/2) -- (0,\bh/2);
      % NB: Computer Modern caps its design size (~25pt) -- a larger \fontsize
      % silently renders no bigger (100pt looked like 26pt). To make the N/S
      % watermark truly fill the bar we SCALE the node instead, which multiplies
      % the glyph geometry and bypasses the cap. \wmscale tunes the height.
      \def\wmscale{1.7}
      \node[figred, opacity=0.10, scale=\wmscale, font=\fontsize{20}{20}\selectfont\sffamily] at (-\bw/4,0) {N};
      \node[ink,    opacity=0.10, scale=\wmscale, font=\fontsize{20}{20}\selectfont\sffamily] at ( \bw/4,0) {S};
    \end{scope}
    % Field arrows on top: two rows (\nrow) symmetric about center, each with \np
    % arrows. A basepoint is a lattice site (bx, by) in the POSITION frame
    % (posang); the arrow there points along the VECTOR frame (vecang). rows are at
    % by = +-\rowsep/2.
    \foreach \row in {1,...,\nrow}{
      \pgfmathsetmacro{\by}{(\row-(\nrow+1)/2)*\rowsep}   % centered rows
      \foreach \k in {1,...,\np}{
        \pgfmathsetmacro{\bx}{-\bw/2 + \k*\bw/(\np+1)}    % site x along the bar
        \pgfmathsetmacro{\px}{#1 + \bx*cos(#3) - \by*sin(#3)}
        \pgfmathsetmacro{\py}{#2 + \bx*sin(#3) + \by*cos(#3)}
        \pgfmathsetmacro{\dx}{-cos(#4)}                    % dir (N->S = -x) by vecang
        \pgfmathsetmacro{\dy}{-sin(#4)}
        \draw[mvec] ({\px-\alen*\dx},{\py-\alen*\dy}) -- ({\px+\alen*\dx},{\py+\alen*\dy});
      }
    }
  }

  % r vector + rotation-axis dot in a panel's POSITION frame. r points from the
  % axis dot to one field basepoint (so it visibly picks out a position). A small
  % gap keeps the arrowhead from touching the bar (they looked "attached"). The
  % blue axis dot is drawn last so it sits above the r arrow.
  \def\rlen{20}   % r length
  \def\rgap{2}    % gap between r's head and the bar edge
  \newcommand{\raxis}[3]{% {cx}{cy}{posang}
    \begin{scope}[shift={(#1,#2)}, rotate=#3]
      \draw[rvec] (0,{-\bh/2-\rlen}) -- (0,{-\bh/2-\rgap});
      \fill[figblue] (0,{-\bh/2-\rlen}) circle[radius=1.1];
    \end{scope}
  }

  % --- Panel centers ----------------------------------------------------------
  \def\colL{0}\def\colR{116}    % left / right column x
  \def\rowT{0}\def\rowB{-82}    % top / bottom row y (extra gap for arrow labels)
  \def\rot{30}                  % the rotation R (degrees)

  % ===== TL: start m. Positions flat, vectors flat.
  \magnet{\colL}{\rowT}{0}{0}
  \raxis{\colL}{\rowT}{0}

  % ===== TR: rotate r. Positions (body + basepoints + r) tilt by R; vectors flat.
  \magnet{\colR}{\rowT}{\rot}{0}
  \raxis{\colR}{\rowT}{\rot}

  % ===== BL: rotate m. Positions flat; vectors tilt by R.
  \magnet{\colL}{\rowB}{0}{\rot}
  \raxis{\colL}{\rowB}{0}

  % ===== BR: correct. Positions AND vectors both tilt by R.
  \magnet{\colR}{\rowB}{\rot}{\rot}
  \raxis{\colR}{\rowB}{\rot}

  % --- Commutative-diagram arrows; all flows reach BR. Every arrow labeled
  % "rotate r" / "rotate m", sloped along the line with the line as baseline
  % (anchor=south via `above`), no background. ------------------------------
  % NB: a node written AFTER the final coordinate (\draw ... -- (end) node{..})
  % attaches at the END point, not the middle -- that put every arrow label at
  % its arrowhead. `midway` forces the geometric center of the segment. `sloped`
  % makes the text run parallel to the line, `above` sets the line as its
  % baseline. No fill => no background.
  % NB: use an ABSOLUTE \fontsize, not \small/\normalsize. The preamble pins
  % every picture to \figfontsize (13.5pt), which overrides relative sizes -- so
  % \normalsize here did nothing. An explicit \fontsize wins.
  % inner sep is the gap between the label and its line; bumped 2pt -> 3pt (1.5x).
  \def\gap{9}
  \tikzset{albl/.style = {midway, sloped, above, font=\fontsize{18}{18}\selectfont, muted, inner sep=3pt}}
  % top edge: TL -> TR (rotate r)
  \draw[tform] ({\colL+\bw/2+\gap},\rowT) -- ({\colR-\bw/2-\gap},\rowT)
  node[albl] {rotate $\boldsymbol{x}$};
  % left edge: TL -> BL (rotate m)
  \draw[tform] (\colL,{\rowT-\bh/2-\rlen-6}) -- (\colL,{\rowB+\bh/2+6})
  node[albl] {rotate $\boldsymbol{m}$};
  % bottom edge: BL -> BR (rotate r)
  \draw[tform] ({\colL+\bw/2+\gap},\rowB) -- ({\colR-\bw/2-\gap},\rowB)
  node[albl] {rotate $\boldsymbol{x}$};
  % right edge: TR -> BR (rotate m). Extended at the TR (top) end by \rext (head
  % fixed at BR); label centered on the (extended) arrow via albl's midway.
  \def\rext{14}
  \draw[tform] (\colR,{\rowT-\bh/2-\rlen-6+\rext}) -- (\colR,{\rowB+\bh/2+6})
  node[albl] {rotate $\boldsymbol{m}$};

  % diagonal: TL -> BR, the correct (composite) transform. Ink + heavier so it
  % reads as the composite. Label centered on the line (midway), sloped, no bg.
  \def\dsx{\colL+\bw/2-1}\def\dsy{\rowT-\bh/2-6}   % start (TL lower-right)
  \def\dex{\colR-\bw/2+1}\def\dey{\rowB+\bh/2+6}   % end   (BR upper-left)
  \draw[line width=1.3pt, ink, -{Stealth[length=2.8mm, width=2.2mm, round]}, line cap=round]
  ({\dsx},{\dsy}) -- ({\dex},{\dey})
  node[midway, sloped, above, font=\fontsize{18}{18}\selectfont, ink, inner sep=3.75pt]
  {rotate $\boldsymbol{x}$ and $\boldsymbol{m}$};

  % --- Panel labels: field label (m / Rm) sits at the OUTER-TOP corner of each
  % magnet; position label (r / Rr) sits by that panel's axis dot. Both are
  % pushed to the column's outer side (left column -> left, right column -> right)
  % so they never fall on the connecting arrows. Rm/Rr, not compositions. -----
  \tikzset{
    flbl/.style = {font=\fontsize{19}{19}\selectfont, figred},
    plbl/.style = {font=\fontsize{19}{19}\selectfont, ink},
  }
  % field labels: outer-top corner of each magnet, rotated with the body so they
  % ride the tilted magnets' top edges.
  \node[flbl, anchor=south east] at ({\colL-\bw/2-3},{\rowT+\bh/2}) {$\boldsymbol{m}$};
  % TR's m and BR's Rm: NOT in the rotated scope -- placed at fixed page positions
  % on TL's / BL's label lines, just right of the tilted magnet's right wall
  % (outer x ~ \colR + (\bw/2)cos(rot) + (\bh/2)sin(rot)). \rlblgap is the gap from
  % that wall; small so the labels sit close to their magnets.
  \def\rlblgap{1}
  \pgfmathsetmacro{\rlblx}{\colR + (\bw/2)*cos(\rot) + (\bh/2)*sin(\rot) + \rlblgap}
  \node[flbl, anchor=west] at ({\rlblx},{\rowB+\bh/2+3}) {$\boldsymbol{m}'$};
  % position labels: centered along the r arrow (its midpoint is at
  % y = -\bh/2 - (\rlen+\rgap)/2), offset to the outer side so it clears the line.
  \def\rmid{-\bh/2-(\rlen+\rgap)/2}
  \node[plbl, anchor=east] at ({\colL-4},{\rowT+\rmid}) {$\boldsymbol{x}$};
  \begin{scope}[shift={(\colR,\rowB)}, rotate=\rot]
    \node[plbl, anchor=west] at ({4},{\rmid}) {$R\boldsymbol{x}$};
  \end{scope}
\end{tikzpicture}
\end{document}
